\documentclass[border=12pt]{standalone}
\usepackage[UTF8,fontset=none]{ctex}
\IfFontExistsTF{Source Han Serif SC}{
  \setmainfont[BoldFont={Source Han Serif SC Bold}]{Source Han Serif SC}
  \setsansfont[BoldFont={Source Han Serif SC Bold}]{Source Han Serif SC}
  \setCJKmainfont[BoldFont={Source Han Serif SC Bold}]{Source Han Serif SC}
  \setCJKsansfont[BoldFont={Source Han Serif SC Bold}]{Source Han Serif SC}
}{
  \PackageWarningNoLine{tikz-infographic}{Source Han Serif SC not found; using Fandol}
  \setCJKmainfont[BoldFont={FandolSong-Bold}]{FandolSong-Regular}
  \setCJKsansfont[BoldFont={FandolHei-Bold}]{FandolHei-Regular}
}
\IfFontExistsTF{Sarasa Mono SC}{
  \setmonofont{Sarasa Mono SC}
  \setCJKmonofont{Sarasa Mono SC}
}{
  \PackageWarningNoLine{tikz-infographic}{Sarasa Mono SC not found; using TeX defaults}
  \setmonofont{Latin Modern Mono}
  \setCJKmonofont{FandolFang-Regular}
}
\usepackage{xcolor}
\usepackage{tikz}
\usetikzlibrary{arrows.meta,backgrounds,calc,fit,positioning,shapes.geometric}
\pgfdeclarelayer{edge layer}
\pgfdeclarelayer{bg layer}
\pgfsetlayers{bg layer,background,edge layer,main}

% ---- Palette ----
\definecolor{Paper}{HTML}{F7F4EE}
\definecolor{Ink}{HTML}{17212B}
\definecolor{Slate}{HTML}{5D6873}
\definecolor{Fog}{HTML}{D9E1E3}
\definecolor{Ocean}{HTML}{2E6F8F}
\definecolor{OceanLight}{HTML}{E2EEF4}
\definecolor{Teal}{HTML}{2A9D8F}
\definecolor{Mint}{HTML}{DDF2EC}
\definecolor{Coral}{HTML}{E76F51}
\definecolor{CoralLight}{HTML}{FBE8E2}
\definecolor{Gold}{HTML}{E9C46A}
\definecolor{GoldLight}{HTML}{FBF1D5}
\definecolor{Violet}{HTML}{7B6FAA}
\definecolor{VioletLight}{HTML}{EAE6F5}
\definecolor{Magma}{HTML}{9F2B68}
\definecolor{Viridian}{HTML}{3C8D6E}

% ---- Styles ----
\tikzset{
  every node/.style={font=\rmfamily,inner sep=0pt,outer sep=0pt},
  panel kicker/.style={anchor=north west,font=\rmfamily\scriptsize\bfseries,text=Teal},
  panel title/.style={anchor=north west,font=\rmfamily\Large\bfseries,text=Ink},
  panel note/.style={anchor=north west,font=\rmfamily\scriptsize,text=Slate,align=left},
  panel bg/.style={draw=Fog!70,fill=white,rounded corners=3mm,inner sep=4mm,line width=.6pt},
  % stack
  stack item/.style={
    rounded corners=2.2mm,inner sep=2.5mm,align=left,
    font=\rmfamily\footnotesize,text=Ink,text width=48mm
  },
  stack primary/.style={stack item,draw=Ocean,fill=Ocean,text=white,font=\rmfamily\bfseries\footnotesize},
  stack secondary/.style={stack item,draw=Ocean!35,fill=OceanLight},
  stack data/.style={stack item,draw=Gold,fill=GoldLight},
  stack output/.style={stack item,draw=Magma,fill=Magma!8,text=Magma!90!black},
  % scene card
  scene card/.style={
    draw=Fog,fill=white,rounded corners=1.6mm,line width=.6pt,
    inner sep=2mm,text width=49mm,
    font=\rmfamily\scriptsize,text=Ink,align=left,anchor=north west
  },
  % pipeline
  pipe node/.style={
    rounded corners=2mm,inner sep=2.5mm,align=center,
    font=\rmfamily\footnotesize\bfseries,text=Ink,text width=35mm,minimum height=11mm
  },
  pipe src/.style={pipe node,draw=Violet,fill=VioletLight,text=Violet!90!black},
  pipe build/.style={pipe node,draw=Ocean,fill=OceanLight,text=Ocean!90!black},
  pipe browser/.style={pipe node,draw=Teal,fill=Mint,text=Viridian},
  pipe output/.style={pipe node,draw=Magma,fill=Magma!10,text=Magma!90!black},
  % gate
  gate node/.style={
    diamond,draw=Coral,fill=CoralLight,inner sep=1.5mm,
    font=\rmfamily\scriptsize\bfseries,text=Coral!90!black,align=center,text width=22mm
  },
  % arrows
  primary flow/.style={
    -{Latex[length=2.4mm,width=1.7mm]},draw=Ocean,
    line width=1.05pt,rounded corners=2mm
  },
  data flow/.style={
    -{Latex[length=2.1mm,width=1.5mm]},draw=Gold!70!black,
    line width=.85pt,rounded corners=1.5mm
  },
  verify flow/.style={
    -{Latex[length=2.1mm,width=1.5mm]},draw=Coral,
    line width=.85pt,rounded corners=1.5mm,dashed
  },
  edge label/.style={
    fill=Paper,inner sep=1.5pt,font=\rmfamily\scriptsize,text=Slate
  },
  % metric
  metric/.style={
    rounded corners=1.6mm,inner sep=2.5mm,align=left,
    font=\rmfamily\bfseries,text=white,text width=36mm,minimum height=14mm
  },
  % badge
  badge/.style={
    rounded corners=1.5mm,inner sep=2mm,
    font=\rmfamily\scriptsize\bfseries,align=center
  },
}

\begin{document}
\begin{tikzpicture}[x=1mm,y=1mm]

  % ---- Page background ----
  \begin{pgfonlayer}{bg layer}
    \fill[Paper] (0,8) rectangle (176,-330);
  \end{pgfonlayer}

  % ============================================================
  % HEADER
  % ============================================================
  \node[panel kicker] at (4,4) {MOSAIC-EXAMPLES \,$\cdot$\, 浏览器内数据库驱动可视化};
  \node[anchor=north west,font=\rmfamily\Huge\bfseries,text=Ink]
    at (4,-1) {百万行数据的透明画布};
  \node[anchor=north west,font=\rmfamily\small,text=Slate,text width=168mm,align=left]
    at (4,-11) {Mosaic + vgplot + DuckDB-WASM 全本地渲染管线；十二场景透明 PNG 参考集，像素级校验确保每张图都有真实数据墨迹与彩色数据标记。};

  % Metrics
  \node[metric,fill=Ocean,anchor=north west] (m1) at (4,-20) {
    {\Large 12\quad}{\footnotesize 个场景}\\[-1mm]
    {\tiny\color{white!85}分析问题全覆盖}
  };
  \node[metric,fill=Teal,anchor=north west] (m2) at (47,-20) {
    {\Large 4.06\quad}{\footnotesize M 行}\\[-1mm]
    {\tiny\color{white!85}总合成数据量}
  };
  \node[metric,fill=Gold,anchor=north west] (m3) at (90,-20) {
    {\Large 0\quad}{\footnotesize 外部请求}\\[-1mm]
    {\tiny\color{white!85}完全离线运行}
  };
  \node[metric,fill=Violet,anchor=north west] (m4) at (133,-20) {
    {\Large 3\quad}{\footnotesize 重验证门}\\[-1mm]
    {\tiny\color{white!85}透明 / 墨迹 / 饱和}
  };

  % ============================================================
  % PANEL 1: 技术栈
  % ============================================================
  \node[panel kicker] (k1) at (4,-38) {01 \,$\cdot$\, 技术栈};
  \node[panel title] (t1) at (4,-42) {三层架构：数据库 \,$\cdot$\, 协调器 \,$\cdot$\, 声明式绘图 API};

  % Column 1: DuckDB-WASM
  \node[stack primary,anchor=north west] (dl1) at (4,-52) {DuckDB-WASM};
  \node[stack secondary,anchor=north west] (dl2) at (4,-63) {
    \textbf{本地列式数据库}\\[.6mm]
    \footnotesize 纯 WASM 运行，无需服务端；\\
    支持 mvp / eh 双 bundle 自动选择
  };
  \node[stack data,anchor=north west] (dl3) at (4,-81) {
    \textbf{合成数据生成}\\[.6mm]
    \footnotesize SQL \texttt{range()} + 三角函数\\
    生成确定性测试数据集
  };

  % Column 2: Mosaic
  \node[stack primary,anchor=north west] (ml1) at (62,-52) {Mosaic Coordinator};
  \node[stack secondary,anchor=north west] (ml2) at (62,-63) {
    \textbf{查询协调 + 缓存}\\[.6mm]
    \footnotesize 将可视化需求编译为 SQL；\\
    结果缓存与查询合并，减少重算
  };
  \node[stack data,anchor=north west] (ml3) at (62,-81) {
    \textbf{DuckDBWASMConnector}\\[.6mm]
    \footnotesize 桥接 vgplot 与本地 DuckDB；\\
    异步管道，支持交叉筛选
  };

  % Column 3: vgplot
  \node[stack primary,anchor=north west] (vl1) at (120,-52) {vgplot API};
  \node[stack secondary,anchor=north west] (vl2) at (120,-63) {
    \textbf{声明式绘图}\\[.6mm]
    \footnotesize \texttt{vg.plot()} + 标记组合；\\
    heatmap / denseLine / hexbin / cell
  };
  \node[stack output,anchor=north west] (vl3) at (120,-81) {
    \textbf{交互原语}\\[.6mm]
    \footnotesize crossfilter 选择 \,$\cdot$\, panZoom\\
    intervalXY 刷选 \,$\cdot$\, 联动视图
  };

  % Arrows
  \begin{pgfonlayer}{edge layer}
    \draw[data flow] (dl1.east) -- (ml1.west)
      node[edge label,midway,above=.3mm] {SQL 查询};
    \draw[data flow] (ml1.east) -- (vl1.west)
      node[edge label,midway,above=.3mm] {聚合结果};
    \draw[primary flow] (dl3.east) -- (ml3.west)
      node[edge label,midway,above=.3mm] {连接器};
    \draw[primary flow] (ml3.east) -- (vl3.west)
      node[edge label,midway,above=.3mm] {API 上下文};
  \end{pgfonlayer}

  % Panel 1 bg
  \begin{pgfonlayer}{bg layer}
    \node[panel bg,fit={(k1.north west) (172,-96)},inner sep=0mm] (p1bg) {};
  \end{pgfonlayer}

  % ============================================================
  % PANEL 2: 场景目录
  % ============================================================
  \node[panel kicker] (k2) at (4,-104) {02 \,$\cdot$\, 场景目录};
  \node[panel title] (t2) at (4,-108) {十二种分析模式，全部透明背景 PNG 输出};
  \node[panel note,text width=168mm] (n2) at (4,-114) {
    每个场景对应一类常见分析问题：从百万点密度到信号监控，从联动视图到运营看板。
    数据量从 18 万到 100 万行不等，覆盖不同量级与可视化家族。
  };

  % Row 1
  \node[scene card,fill=Magma!8] (c1) at (4,-125) {
    {\ttfamily\tiny\bfseries\color{Magma} density}\hfill{\tiny\color{Slate}1,000,000 行}\\[.5mm]
    {\bfseries\scriptsize 百万点密度热图}\\[.2mm]
    {\scriptsize\color{Slate}屏幕分辨率下，一百万点集中在哪？}
  };
  \node[scene card,fill=Ocean!8] (c2) at (62,-125) {
    {\ttfamily\tiny\bfseries\color{Ocean} signals}\hfill{\tiny\color{Slate}400,000 行}\\[.5mm]
    {\bfseries\scriptsize 密集时序信号图集}\\[.2mm]
    {\scriptsize\color{Slate}四条长时序如何保持可探索性？}
  };
  \node[scene card,fill=Teal!8] (c3) at (120,-125) {
    {\ttfamily\tiny\bfseries\color{Teal} linked}\hfill{\tiny\color{Slate}320,000 行}\\[.5mm]
    {\bfseries\scriptsize 联动操作节律}\\[.2mm]
    {\scriptsize\color{Slate}空间刷选如何过滤小时剖面？}
  };

  % Row 2
  \node[scene card,fill=Gold!12] (c4) at (4,-144) {
    {\ttfamily\tiny\bfseries\color{Gold!70!black} distribution}\hfill{\tiny\color{Slate}360,000 行}\\[.5mm]
    {\bfseries\scriptsize 数据库驱动分布对比}\\[.2mm]
    {\scriptsize\color{Slate}四个队列在密集测量上有何差异？}
  };
  \node[scene card,fill=Violet!8] (c5) at (62,-144) {
    {\ttfamily\tiny\bfseries\color{Violet} regression}\hfill{\tiny\color{Slate}480,000 行}\\[.5mm]
    {\bfseries\scriptsize 密集关系探索}\\[.2mm]
    {\scriptsize\color{Slate}密集采样后还剩什么非线性关系？}
  };
  \node[scene card,fill=Coral!8] (c6) at (120,-144) {
    {\ttfamily\tiny\bfseries\color{Coral} retention}\hfill{\tiny\color{Slate}320,000 行}\\[.5mm]
    {\bfseries\scriptsize 队列留存图集}\\[.2mm]
    {\scriptsize\color{Slate}八个队列 24 期内如何留存？}
  };

  % Row 3
  \node[scene card,fill=Viridian!8] (c7) at (4,-163) {
    {\ttfamily\tiny\bfseries\color{Viridian} quality-matrix}\hfill{\tiny\color{Slate}280,000 行}\\[.5mm]
    {\bfseries\scriptsize 质量风险矩阵}\\[.2mm]
    {\scriptsize\color{Slate}哪些阶段与子系统累积质量风险？}
  };
  \node[scene card,fill=Ocean!8] (c8) at (62,-163) {
    {\ttfamily\tiny\bfseries\color{Ocean} rankings}\hfill{\tiny\color{Slate}240,000 行}\\[.5mm]
    {\bfseries\scriptsize 区域绩效排名}\\[.2mm]
    {\scriptsize\color{Slate}区域评分与不确定性如何比较？}
  };
  \node[scene card,fill=Magma!8] (c9) at (120,-163) {
    {\ttfamily\tiny\bfseries\color{Magma} anomalies}\hfill{\tiny\color{Slate}420,000 行}\\[.5mm]
    {\bfseries\scriptsize 异常事件信号监控}\\[.2mm]
    {\scriptsize\color{Slate}稀疏尖峰如何在密集信号中显现？}
  };

  % Row 4
  \node[scene card,fill=Teal!8] (c10) at (4,-182) {
    {\ttfamily\tiny\bfseries\color{Teal} vector-field}\hfill{\tiny\color{Slate}180,000 行}\\[.5mm]
    {\bfseries\scriptsize 采样流场向量}\\[.2mm]
    {\scriptsize\color{Slate}方向流如何随采样网格变化？}
  };
  \node[scene card,fill=Violet!8] (c11) at (62,-182) {
    {\ttfamily\tiny\bfseries\color{Violet} small-multiples}\hfill{\tiny\color{Slate}360,000 行}\\[.5mm]
    {\bfseries\scriptsize 服务行为小倍数}\\[.2mm]
    {\scriptsize\color{Slate}同指标在四个服务中如何表现？}
  };
  \node[scene card,fill=Gold!12] (c12) at (120,-182) {
    {\ttfamily\tiny\bfseries\color{Gold!70!black} operations}\hfill{\tiny\color{Slate}520,000 行}\\[.5mm]
    {\bfseries\scriptsize 联动运营驾驶舱}\\[.2mm]
    {\scriptsize\color{Slate}负载/延迟/错误/队列如何对齐？}
  };

  % Bar chart section
  \node[panel kicker] (bk) at (4,-205) {数据规模一览};
  \node[panel note,text width=168mm] (bn) at (4,-209) {
    总计 \textbf{4,060,000} 行合成数据，全部由 SQL 在本地 DuckDB 内生成，无网络请求、无外部依赖。
  };

  % Inline bar chart (bars go downward from baseline at y=-214)
  \begin{scope}[shift={(8,-214)}]
    \def\barW{10.5}
    \def\gap{2.7}
    \def\maxH{50}
    % baseline (top)
    \draw[draw=Fog,line width=.6pt] (0,0) -- (158,0);
    % y-axis labels (top to bottom: 1M at top, 0 at bottom)
    \node[font=\tiny,text=Slate,anchor=east] at (-2, 0) {1M};
    \node[font=\tiny,text=Slate,anchor=east] at (-2, -25) {500K};
    \node[font=\tiny,text=Slate,anchor=east] at (-2, -\maxH) {0};
    \node[font=\tiny,text=Slate,anchor=east,rotate=90] at (-7, -25) {行数};
    % bars (draw downward: y=0 to y=-h)
    \foreach \i/\h/\col/\lab in {
      0/50/Magma/density, 1/20/Ocean/signals, 2/16/Teal/linked,
      3/18/Gold/dist, 4/24/Violet/regr, 5/16/Coral/reten,
      6/14/Viridian/qmatrix, 7/12/Ocean/rank, 8/21/Magma/anom,
      9/9/Teal/vector, 10/18/Violet/small, 11/26/Gold/ops
    } {
      \pgfmathsetmacro\xpos{\i * (\barW + \gap)}
      \draw[fill=\col!45,draw=\col!60,line width=.4pt,rounded corners=.6mm]
        (\xpos, 0) rectangle (\xpos + \barW, -\h);
      \node[font=\ttfamily\tiny,text=Slate,anchor=north,rotate=45,inner sep=.5pt]
        at (\xpos + \barW/2, -\maxH - 1.5) {\lab};
    }
  \end{scope}

  % Panel 2 bg
  \begin{pgfonlayer}{bg layer}
    \node[panel bg,fit={(k2.north west) (172,-272)},inner sep=0mm] (p2bg) {};
  \end{pgfonlayer}

  % ============================================================
  % PANEL 3: 渲染与验证管线
  % ============================================================
  \node[panel kicker] (k3) at (4,-280) {03 \,$\cdot$\, 渲染与验证管线};
  \node[panel title] (t3) at (4,-290) {从源码到透明 PNG：单命令全量验证};

  % Pipeline top row
  \node[pipe src,anchor=north west] (src) at (4,-300) {
    \textbf{源码}\\[.6mm]
    \footnotesize\texttt{src/main.ts}\\
    \footnotesize\texttt{catalog.json}
  };
  \node[pipe build,anchor=north west] (build) at (48,-300) {
    \textbf{Vite 构建}\\[.6mm]
    \footnotesize\texttt{vite build} $\to$ \texttt{dist/}
  };
  \node[pipe browser,anchor=north west] (pw) at (92,-300) {
    \textbf{Playwright 浏览器}\\[.6mm]
    \footnotesize 无头 Chromium + Chrome
  };
  \node[pipe output,anchor=north west] (out) at (136,-300) {
    \textbf{透明 PNG 输出}\\[.6mm]
    \footnotesize\texttt{out/*-transparent.png}
  };

  % Preview server below build
  \node[pipe build,anchor=north west,fill=white,draw=Ocean!40,text=Ocean!70!black] (preview) at (48,-316) {
    \textbf{预览服务器}\\[.6mm]
    \footnotesize\texttt{vite preview :41732}
  };

  % Arrows (top row)
  \begin{pgfonlayer}{edge layer}
    \draw[primary flow] (src.east) -- (build.west)
      node[edge label,midway,above=.3mm] {tsc + vite};
    \draw[primary flow] (build.east) -- (pw.west)
      node[edge label,midway,above=.3mm] {静态页面};
    \draw[primary flow] (pw.east) -- (out.west)
      node[edge label,midway,above=.3mm] {截图};
    \draw[data flow] (build.south) -- (preview.north)
      node[edge label,midway,right=1mm,xshift=.5mm] {serve};
  \end{pgfonlayer}

  % Validation gates
  \node[panel kicker] (vk) at (4,-334) {三重像素级验证门};

  \node[gate node,anchor=north west] (g1) at (22,-342) {透明度\\校验};
  \node[gate node,anchor=north west] (g2) at (78,-342) {墨迹像素\\校验};
  \node[gate node,anchor=north west] (g3) at (134,-342) {色彩饱和\\度校验};

  % Gate descriptions
  \node[font=\rmfamily\scriptsize,text=Slate,text width=38mm,align=center,anchor=north] (gd1) at (22,-358) {
    透明像素 > 18\%\\确保背景真正透明
  };
  \node[font=\rmfamily\scriptsize,text=Slate,text width=38mm,align=center,anchor=north] (gd2) at (78,-358) {
    可见像素 > 8,000\\确保图形非空
  };
  \node[font=\rmfamily\scriptsize,text=Slate,text width=38mm,align=center,anchor=north] (gd3) at (134,-358) {
    饱和像素 > 1,200\\确保彩色数据标记
  };

  % Verify flow arrows (from output down, then to three gates)
  \begin{pgfonlayer}{edge layer}
    \draw[verify flow] (out.south) -- ++(0,-26mm) coordinate (vbranch);
    \draw[verify flow] (vbranch) -| (g1.north);
    \draw[verify flow] (vbranch) -| (g2.north)
      node[edge label,pos=0.5,above=.3mm] {PNG 像素分析};
    \draw[verify flow] (vbranch) -| (g3.north);
  \end{pgfonlayer}

  % Badge row
  \node[badge,draw=Coral!60,fill=CoralLight,text=Coral!90!black,anchor=north west] (b1) at (4,-372)
    {阻断所有外部请求};
  \node[badge,draw=Teal!60,fill=Mint,text=Viridian,anchor=north west] (b2) at (42,-372)
    {指针交互测试};
  \node[badge,draw=Ocean!60,fill=OceanLight,text=Ocean!90!black,anchor=north west] (b3) at (80,-372)
    {\texttt{\_\_mosaicDemo.ready} 就绪标志};

  % Panel 3 bg
  \begin{pgfonlayer}{bg layer}
    \node[panel bg,fit={(k3.north west) (172,-380)},inner sep=0mm] (p3bg) {};
  \end{pgfonlayer}

  % ============================================================
  % FOOTER
  % ============================================================
  \node[anchor=north west,font=\rmfamily\scriptsize,text=Slate,text width=168mm]
    at (4,-388) {
    \textbf{项目路径} \texttt{~/projects/plot/mosaic-examples}
    \,$\cdot$\, \textbf{构建} \texttt{just build}
    \,$\cdot$\, \textbf{渲染} \texttt{just render}
    \,$\cdot$\, \textbf{验证} \texttt{just test}
  };
  \node[anchor=north west,font=\rmfamily\scriptsize,text=Slate,text width=168mm]
    at (4,-393) {
    数据来源：\texttt{catalog.json} \,$\cdot$\, \texttt{src/main.ts} \,$\cdot$\, \texttt{tools/validate.mjs} \,$\cdot$\, \texttt{package.json}
    \,$\mid$\, 生成日期：2026-08-24
  };

  % Expand page background to fit
  \begin{pgfonlayer}{bg layer}
    \fill[Paper] (0,8) rectangle (176,-400);
  \end{pgfonlayer}

\end{tikzpicture}
\end{document}
